\paragraph{$\bm{n}$-$\bm{n}$}

\paragraph{m-m}\label{m-m}

\[
\begin{aligned}
\LinearMM &= \SectionIntegral{
     [\bm{r}^\times][\bm{r}^\times]^{\mathrm{t}} 
}
= \SectionIntegral{ \left[\begin{matrix}
y^{2} + z^{2} & 0 & 0\\
0 & z^{2} & - z y \\
0 & - z y & \,\,y^{2}
\end{matrix}\right]
}
% &= \int (\bm{r}\cdot\bm{r})\bm{1} - \bm{r}\otimes \bm{r}\, \mathrm{d}A \\
% &=\left[\delta_{\alpha \beta} \mathbf{1}-\FrameBasis_\alpha \otimes \FrameBasis_\beta\right]\int_{\Omega} \xi_\alpha \xi_\beta \, \mathrm{d}A  \\
\\
% &\triangleq I_0 \FrameBasis \otimes \FrameBasis + I_{\alpha \beta} \FrameBasis_\alpha \otimes \FrameBasis_\beta  \\
\end{aligned}
\]
\textbf{three} independent constants.
% \[
% \begin{aligned}
% I_0    &\triangleq  y^2 + z^2 \\
% I_y    &\triangleq  I_{22} = \int z^2\, \mathrm{d} A \\
% I_z    &\triangleq  I_{33} = \int y^2\, \mathrm{d} A \\
% I_{zy} &\triangleq -I_{23} = \int yz \, \mathrm{d} A
% \end{aligned}
% \]
Note
\[
\LinearMM = \SectionIntegral{\bm{r}\cdot\bm{r}\bm{1} - \bm{r}\otimes\bm{r}}
\]

\begin{itemize}
    \item The \emph{principal axes} of the cross section are defined as the coordinate system which diagonalizes \(\LinearMM\). Consequently, in this coordinate system the resultant moments \(M_{\beta}\) are often decoupled. 
\end{itemize}

\paragraph{m-m Wagner}

\[
\tau r^{2} (\bm{r} {\times} \FrameBasis \otimes \FrameBasis +  \FrameBasis \otimes \bm{r} {\times}\FrameBasis)
\]

\paragraph{w-w}\label{w-w}

\begin{equation}
\LinearWW = \SectionIntegral{\bm{\varphi}\otimes\bm{\varphi}}
= \SectionIntegral{ \begin{bmatrix}
\varphi^2  & \varphi \psi_y & \varphi \psi_z \\
           & \psi_y^2 & \psi_z \psi_y \\
           &          & \psi_z^2
\end{bmatrix}}
= \begin{bmatrix}
J_{w} & 0 & 0 \\
0 & C_{\psi} & 0 \\
0 & 0 & C_{\psi}
\end{bmatrix}
\end{equation}
where \textbf{6} independent constants are defined. This is equivalent to the \emph{warping matrix} of \cite{elfatmi2007nonuniform}.

\begin{itemize}
\tightlist
\item
  If \(\bm{\varphi} = \tilde{\varphi} \,\FrameBasis\) then \(J_w\) is the warping
  constant of \cite[p39]{vlasov1961thinwalled} 
  \[
  J_w = \SectionIntegral{ \tilde{\varphi}^2 }
  \]
  Some other names for this constant include \emph{sectorial moment of inertia} \cref{vlasov1961thinwalled}, 
\end{itemize}

\paragraph{\(\bm{v}-\bm{v}\)}

\begin{equation}
\LinearVV = \SectionIntegral{\nabla_{\beta} \bm{\varphi}\otimes\nabla_{\beta} \bm{\varphi}}
= \SectionIntegral{ 
\begin{bmatrix}
\varphi_{,y}^2 & \varphi_{,y}\psi_{y,y} & \varphi_{,y}\psi_{z,y} \\
& \psi_{y,y}^2 & \psi_{y,y}\psi_{z,y} \\
& & \psi_{z,y}^2
\end{bmatrix} +
\begin{bmatrix}
\varphi_{,z}^2 & \varphi_{,z}\psi_{y,z} & \varphi_{,z}\psi_{z,z} \\
& \psi_{y,z}^2 & \psi_{y,z}\psi_{z,z} \\
& & \psi_{z,z}^2
\end{bmatrix}
}
\end{equation}
\textcite{gruttmann2000theory} refer to \(J_{v}\) 
% NOTE: I previously used C_{\alpha} for J_v
with the notation ``\(I_0 + I_T\)''
\[
  J_{v}  = \SectionIntegral{ \tilde{\varphi}_{, y}^2 + \tilde{\varphi}_{, z}^2} = J_S - J
\]

\paragraph{Coupling in $\bm{n}-\bm{m}$}\label{n-m}

\[
\LinearNM
\triangleq \SectionIntegral{ -[\bm{r}^\times] }
= \SectionIntegral{ \begin{bmatrix}
 0 & z &-y \\
-z & 0 & 0 \\
 y & 0 & 0
\end{bmatrix} \, \mathrm{d} A
 = \begin{bmatrix}
     0 & Q_{y} & -Q_{z} \\
-Q_{y} & 0 & 0 \\
 Q_{z} & 0 & 0
\end{bmatrix}
}
\]
where \textbf{two} independent constants known as the \emph{first moments of area} have been defined.

\begin{itemize}
\item
  These are exactly an area-scaling of the centroid coordinates \(\CentroidLocation\) defined in \cref{eq:define-centroid}, that is,  
  \[
  \bm{A}_r \equiv A \CentroidLocation^{\times}
  \]
\item
  If the reference curve passes through the centroid as
  in \cite{simo1991geometricallyexact} then \(\bm{A}_r = \bm{0}\)
\end{itemize}

\paragraph{Coupling in $\bm{n}-\bm{w}$}\label{n-w}

\[
\SectionNW \triangleq \SectionIntegral{ \FrameBasis \otimes\bm{\varphi} 
}
= \SectionIntegral{\begin{bmatrix}
\varphi & \psi_y & \psi_z \\
0 & 0 & 0 \\
0 & 0 & 0 \\
\end{bmatrix}
}
% = \begin{bmatrix}
% J_{nw} & S_{ny} & S_{nz} \\
% 0 & 0 & 0 \\
% 0 & 0 & 0
% \end{bmatrix}
\] 
where \textbf{three} independent constants have been defined and
\[
\SectionNV \triangleq \SectionIntegral{ \FrameBasis_{\beta}\otimes \nabla_{\beta}\bm{\varphi}
}
= \SectionIntegral{\begin{bmatrix}
0 & 0 & 0 \\
\varphi_{,y} & \psi_{y,y} & \psi_{z,y} \\
\varphi_{,z} & \psi_{y,z} & \psi_{z,z} \\
\end{bmatrix}
}
= \begin{bmatrix}
0 &     0 & 0 \\
J_{ny} & S_{nyy} & S_{nzy} \\
J_{nz} & S_{nyz} & S_{nzz}
\end{bmatrix}
\] 
assuming \(\bm{\varphi}' = \bm{0}\).

\begin{itemize}
\item If \(\varphi\) is defined by a pure Neumann problem of the Laplacian, then the condition 
\[
\SectionIntegral{\varphi} = 0
\]
is most likely satisfied by \(\varphi\). This is generally the case for St. Venant warping.
\item
  If \(\varphi\) is referred to the \emph{shear
  center} (e.g., \cite{gruttmann1998geometrical}) 
  \begin{equation*}
  J_{ny} = \int \tilde{\varphi}_{,y} = -(\tilde{z} - \CentroidLocationZ) A
  \qquad\text{and}\qquad
  J_{nz} = \int \tilde{\varphi}_{,z} =  (\tilde{y} - \CentroidLocationY) A
  \end{equation*}
\item
  If \(\varphi\) is referred to the \emph{centroid}
  (i.e.~\(\varphi = \hat{\varphi}\) like \cite{gruttmann2000theory}) then:
  \[
  J_{ny} = \int \hat{\varphi}_{,y} =    \CentroidLocationZ A \\
  \qquad\text{and}\qquad
  J_{nz} = \int \hat{\varphi}_{,z} =  - \CentroidLocationY A
  \]
  where \(\CentroidLocationZ\) and \(\CentroidLocationY\) are the coordinates of the centroid.
\end{itemize}

\paragraph{m-w}\label{m-w}
% NOTE: May remove the symbols J_my, etc... if they are always equal to -Iv, then they are redundant.
\[
\LinearMW 
= \SectionIntegral{ \bm{r} \times\FrameBasis \otimes \bm{\varphi} } 
= \SectionIntegral{ \begin{bmatrix}
0 & 0 & 0 \\
~z \varphi & ~z \psi_{y} & ~z \psi_{z} \\
-y \varphi & -y \psi_{y} & -y \psi_z
\end{bmatrix}
}
= \begin{bmatrix}
0 & 0 & 0 \\
J_{my} & S_{myy} & S_{mzy} \\
J_{mz} & S_{myz} & S_{mzz}
\end{bmatrix}
\]
introducing \textbf{two} shear-interacting constants \(J_{my}\) and \(J_{mz}\) for torsion.
\begin{itemize}
\tightlist
\item
  If \(\varphi\) is referred to the centroid 
  \[
  \begin{aligned}
  S_y &= \int z \hat{\varphi} \, \mathrm{d} A =   \hat{I}_{z} \hat{\SectionLocation}_z - \hat{I}_{yz} \hat{y} \\
  S_z &= \int y \hat{\varphi} \, \mathrm{d} A = - \hat{I}_{y} \hat{\SectionLocation}_y + \hat{I}_{yz} \hat{z}
  \end{aligned}
  \]
\end{itemize}
For moment-bishear interaction:
\[
\LinearMV
= \SectionIntegral{ \bm{r} \times\FrameBasis_{\beta} \otimes \nabla_{\beta}\bm{\varphi} } 
= \SectionIntegral{ \begin{bmatrix}
-z \varphi_{, y} + y \varphi_{, z} &
-z \psi_{y, y} + y \psi_{y, z} &
-z \psi_{y, y} + y \psi_{y, z} \\
0 & 0 & 0 \\
0 & 0 & 0 \\
\end{bmatrix}
}
% = \begin{bmatrix}
% J_{mw} & S_{my} & S_{mz} \\
% 0 & 0 & 0 \\
% 0 & 0 & 0 \\
% \end{bmatrix}
\]
Generally, however, \(\LinearMV \equiv - \LinearVV\).
% \[
% \SectionMW 
% = \SectionIntegral{ \bm{r} \times\FrameBasis \otimes \bm{\varphi} } 
% = \SectionIntegral{ \begin{bmatrix}
% 0 & -z \varphi_{, y} + y \varphi_{, z} \\
%  z \varphi & 0 \\
% -y \varphi & 0
% \end{bmatrix}
% }
% = \begin{bmatrix}
% 0 & S_{\alpha} \\
% S_{y} & 0 \\
% S_{z} & 0
% \end{bmatrix}
% \]

